\chapter{結論}
\label{chap:conclusion}

%% ===========================================================
\section{本研究のまとめ}
\label{sec:summary}
本論文では，\upLaTeX による卒業論文・修士論文の執筆を支援するテンプレートについて述べた．
第\ref{chap:writing}章では論文執筆の作法を，第\ref{chap:latex}章および第\ref{chap:float}章では
論文の執筆に必要な\LaTeX の記法を，さまざまなパターンとともに網羅的に示した．
これにより，執筆者は体裁を整える手間を抑え，内容の検討に集中できる．

結論では，このように序論（第\ref{chap:intro}章）で述べた目的に対して，
本研究で何を行い，どの程度達成できたかを簡潔に振り返る．

%% ===========================================================
\section{今後の課題}
\label{sec:future}
研究の限界や今後の課題・展望を述べることで，研究の発展性を示す．
結論が明確でない，あるいは目的の内容と結論が乖離しているものは再考を要する\cite{ipsj-sample}．
本テンプレートについては，研究分野ごとのスタイル整備や，BibTeXによる文献管理の
解説の拡充などが今後の課題として挙げられる．
